% Independent arithmetic audit. Statement IDs reserved with the parent.
% This file uses the result environment with arguments id, status, title.

\subsection{The completion and finite-support reconstruction}

\begin{result}{E.101}{Proved}{Finite support; SECTIONS\_OVER\_Q A.32 and GT\_TOWER\_CHECK C.1.1 confirmed}
Put $\mathcal K=\varprojlim_n\mathbb Q^\times/(\mathbb Q^\times)^n$,
with discrete power quotients and their inverse-limit topology. We have a
continuous group isomorphism, with continuous inverse,
\[
\mathcal K\simeq\{\pm1\}\times
\left\{(a_q)_q\in\prod_q\widehat{\mathbb Z}:
  \text{for every }n,\ a_q=0\pmod n\text{ for almost all }q\right\}.
\]
The topology on the right is defined by simultaneous reduction of all
coordinates modulo $n$, with the resulting direct sum given the discrete
topology. If $(A,B)$ is ordinary at every finite prime, every valuation
$v_q(A),v_q(B)$ is an ordinary integer. If any one of $A,B,A/B$ has finite
valuation support, there is a unique $t\in\mathbb Q\setminus\{0,1\}$ with
$(A,B)=(\kappa(t),\kappa(1-t))$; all its ordinary local points equal $t$.
No local section theorem is used.
\end{result}
\begin{proof}
Unique factorization identifies $\mathbb Q^\times$ with
$\{\pm1\}\times\bigoplus_q\mathbb Z$. At each power quotient the valuation
coordinates have finite support. A compatible collection of these finite
support vectors gives exactly the displayed restricted family, and vice
versa. Both constructions commute with the defining quotient maps, proving
the assertion about topology as well as the group isomorphism.

The local multiplicative Kummer map
$\mathbb Q_p^\times\to\varprojlim_n\mathbb Q_p^\times/(\mathbb Q_p^\times)^n$
is injective. Indeed its valuation has to be divisible by every $n$, and
is zero. The unit group is a product of a finite group and a group
isomorphic to $\mathbb Z_p$; at two one first uses $1+4\mathbb Z_2$.
The intersection of all power subgroups in each factor is trivial.
Ordinarity therefore gives $v_q(A)=v_q(x_q)\in\mathbb Z$, and similarly
for $B$.

If $A$ has finite support, its sign and integer valuations give an actual
rational number $t$, and the displayed isomorphism gives $A=\kappa(t)$.
Local injectivity gives $x_p=t$ for every $p$, in particular $t\ne0,1$.
The quotient $B/\kappa(1-t)$ is locally trivial everywhere. Its local
valuations annihilate every prime coordinate; its remaining possible sign
is detected at $p=3$, since $-1$ is not a square in $\mathbb Q_3$.
Thus $B=\kappa(1-t)$. Interchanging the coordinates handles finite support
of $B$. If $A/B$ has finite support, it equals $\kappa(u)$ for
$u\in\mathbb Q^\times$, and local injectivity gives
$x_p/(1-x_p)=u$. The value $u=-1$ is impossible, and hence
$x_p=u/(1+u)$ for every $p$. The same local-triviality argument recovers
both global coordinates. Uniqueness follows from any local Kummer map.
\end{proof}

% The following two statements belong in Part B when assembling the paper.
\subsection{A unique class-two candidate for the original factorial ghosts}

\begin{result}{B.101}{Proved}{Two exceptional-prime tests select a unique class-two class}
Let $(A,B)\in\mathcal G_{\rm ord}$ satisfy
\[
(v_3(A),v_3(B))=(1,0),\qquad
(v_2(A),v_2(B))=(0,1).
\]
Use the actual integral class-two arithmetic quotient of the tripod,
the tangential splitting at zero, and geometric conjugacy of continuous
sections. There is exactly one class-two section class above $(A,B)$
whose third obstruction vanishes at both of the following pairs of
rational place and coefficient prime:
\[
(q,\ell)=(3,2),\qquad(q,\ell)=(2,3).
\]
Consequently at most one class-two section class above this pair can
be the projection of an integral class-three section. In particular
the twenty-four alternatives in the original A.44 family reduce to
one canonical candidate before the other local or continuous tests
are imposed. The assertion does not say that this candidate lifts.
\end{result}
\begin{proof}
We use the independently checked weight-two calculation
$H^1(G_{\mathbb Q},\widehat{\mathbb Z}(2))=\mathbb Z/24$, generated
by the class of $\eta(g)=(\chi(g)^2-1)/24$, together with the actual
class-two lifting result for ordinary pairs. For fixed literal abelian
cocycles $a,b$, write a section in Hall coordinates as
$y^{b(g)}x^{a(g)}C^{h(g)}s_0(g)$. Its equation is
$\delta h=-a\cup b$. Every other class-two class is represented by
$h+k\eta$, for one $k\in\mathbb Z/24$. Central coboundaries account
for the change of representative. There is no additional abelian
stabilizer: $H^0(G_{\mathbb Q},\widehat{\mathbb Z}(1)^2)=0$,
since each local factor is torsion-free and a nonidentity cyclotomic
scalar annihilates no nonzero vector.

At $q=3$, with coefficient prime two, ordinarity and $v_3(A)=1$
make the local class of $B$ that of the principal unit $1-x_3$,
and hence zero in $H^1(G_{\mathbb Q_3},\mathbb Z_2(1))$.
Thus the literal cocycle $b$ is $\delta u$ for a unique
$u\in\mathbb Z_2(1)$. For tame inertia $T$ and arithmetic Frobenius
$F$, put $\alpha=a(F)$ and $\gamma=h(F)$. We have
\[
a(T)=1,\quad b(T)=0,\quad b(F)=2u,\quad h(T)=u.
\]
The last equality follows from the actual class-two tame relation;
the scalar $q-1$ is not a zero divisor in the integral module, so
there is no extra inertia freedom. Substitution into the actual third
residue formula gives
\[
Z_D=\gamma-3\alpha u,\qquad
Z_F=-8\binom{u+1}{2}.
\]
Here we use the actual cochain and purity comparison of B.22--B.23.
The local obstruction group is
$H^2(G_{\mathbb Q_3},\mathbb Z_2(3))=\mathbb Z_2/8$ in each
coordinate. The second displayed residue is already zero there.
Twisting $h$ by $k\eta$ leaves its inertia unchanged and changes
$\gamma$ by $k\eta(F)=k/3$. Multiplication by $1/3$ is a unit
modulo eight. Exactly one $k\pmod8$ therefore makes both residues
zero.

At $q=2$, with coefficient prime three, the ordinary local class of
$A=1-(1-x_2)$ is a principal two-adic unit and is zero in the
three-primary Kummer group. Write $a=\delta v$, uniquely, and set
$\beta=b(F)$ and $\gamma=h(F)$ in this component. The tame data are
\[
a(T)=0,\quad b(T)=1,\quad a(F)=v,\quad h(T)=-v.
\]
The residue formulas now give
\[
Z_D=3\binom{v+1}{2},\qquad
Z_F=\gamma+\beta v-\frac18.
\]
The local obstruction group in each coordinate is $\mathbb Z_3/3$.
Again the first residue already vanishes, and twisting changes the
second by $k\eta(F)=k/8$, a unit multiple of $k$ modulo three.
Exactly one $k\pmod3$ makes this local obstruction zero.

The two congruences select exactly one $k\pmod{24}$ by the Chinese
remainder theorem. All formulas are insensitive to central coboundaries,
which change the Frobenius values by multiples of $q^2-1$.
Changing aligned cocycle representatives amounts to geometric
conjugacy and preserves the obstruction. This proves existence and
uniqueness for these two tests, and the necessity for a global
class-three lift follows by continuous restriction to the two local
groups. The polynomial substitutions and the resulting bijection
$k\mapsto(k/3\bmod8,k/8\bmod3)$ are checked in
\texttt{code/ghost\_audit\_checks.py}.
\end{proof}

\begin{result}{B.102}{Proved}{The remaining fresh-support test in two scalar sequences}
For an original A.44 pair let $s$ be the canonical class-two candidate
of B.101 and put $n_j=(j+4)!$. Fix a coefficient prime $\ell$.
At each positive $A$-support prime $q_j\ne\ell$, write
$b=\delta u_j$ locally, and use Frobenius coordinates
$\alpha_j=a(F_{q_j})$, $\gamma_j=h(F_{q_j})$.
At each positive $B$-support prime $r_j\ne\ell$, write
$a=\delta v_j$ and put
$\beta_j=b(F_{r_j})$, $\widetilde\gamma_j=h(F_{r_j})$.
The two potentially nonintegral normalized residues are
\[
T_{A,j}=\frac{n_j(\gamma_j-q_j\alpha_j u_j)}{q_j^2-1},
\qquad
T_{B,j}=\frac{n_j(\widetilde\gamma_j+\beta_jv_j)}{r_j^2-1}.
\]
The fresh-support local class-three obstructions all vanish if and
only if these numbers belong to $\mathbb Z_\ell$ for every indicated
index and every $\ell$. After the finitely many old-place residues
have also been tested, the global continuous class-three obstruction
vanishes if and only if, in addition,
$T_{A,j}\to0$ and $T_{B,j}\to0$ in $\mathbb Z_\ell$ for every $\ell$.
The quotient formulas are exact, including coefficient primes two
and three. This is a reduction of B1 to explicit arithmetic conditions
on one section class, not a construction of a prime sequence satisfying
them.
\end{result}
\begin{proof}
We use the strict continuous localization calculation of A.22 and
the literal cochain comparison B.22--B.23, in
which the third obstruction is a class in
$R_\ell/(q^2-1)_qR_\ell$ for each of its two central coordinates,
where $R_\ell$ consists of families tending to zero. Local vanishing
is coordinatewise divisibility by $q^2-1$; global vanishing requires
the divided family itself to lie in $R_\ell$.

At $q=q_j$ the data are
$m=n_j$, $n=0$, $t=n_ju_j$, and $\beta=(q-1)u_j$.
Direct substitution into the two third-residue polynomials gives
\[
\frac{Z_{D,q}}{q^2-1}
 =T_{A,j}+u_j\binom{n_j}{2},\qquad
\frac{Z_{F,q}}{q^2-1}
 =-n_j\binom{u_j+1}{2}.
\]
At $q=r_j$ we instead have $m=0$, $n=n_j$, $t=-n_jv_j$,
and $\alpha=(q-1)v_j$. The exact identities are
\[
\frac{Z_{D,q}}{q^2-1}
 =n_j\binom{v_j+1}{2},\qquad
\frac{Z_{F,q}}{q^2-1}
 =T_{B,j}-\frac{n_j}{24}-v_j\binom{n_j}{2}.
\]
The term $n_j/24$ retains the actual Ihara-action correction; it
must not be discarded at two or three. It is integral because every
original factorial $n_j$ is divisible by twenty-four.

Each $u_j,v_j$ belongs to the integral local Tate module after the
chosen trivialization. The binomial expression $z(z+1)/2$ is integral
on every $\mathbb Z_\ell$, including $\mathbb Z_2$. Since
$n_j\to0$ in every $\mathbb Z_\ell$, all the other terms in the
four displayed quotients tend to zero: the only possible fixed
denominators are two and twenty-four. They are therefore irrelevant
both to divisibility and, after integrality is established, to the
zero-limit condition. This proves the equivalence at fresh support.
At zero-valuation primes the integral class-two tame relation forces
all inertia coordinates to vanish, so both residues are zero. The
old primes two and three contribute finitely many terms, which must
be tested for integrality but do not alter convergence of the family.
A.22 then gives the asserted global equivalence. No local section
theorem enters any of these computations.
\end{proof}

\subsection{An independent reconstruction of the original arithmetic ghost}

\begin{result}{E.102}{Proved}{SECTIONS\_OVER\_Q A.44 confirmed: ordinary infinite-support pairs}
There are pairwise distinct primes $q_j,r_j>3$, for $j\ge0$, such that
\[
A=3\prod_{j\ge0}q_j^{(j+4)!},\qquad
B=-2\prod_{j\ge0}r_j^{(j+4)!}
\]
converge in $\mathcal K$. At every finite prime $p$ they are the Kummer
classes of nonzero $x_p,y_p\in\mathbb Z_p$ satisfying $x_p+y_p=1$.
Their nonzero valuation vectors are $(1,0)$ at three, $(0,1)$ at two,
$((j+4)!,0)$ at $q_j$, and $(0,(j+4)!)$ at $r_j$. The first prime
$q_0$ can take infinitely many values. The result concerns the abelian
completed pair only; it does not assert a section of any nonabelian
quotient. No local section theorem is used.
\end{result}
\begin{proof}
We retain the explicit factorials of A.44, and independently check the
field intersection and the induction on which its assertion depends.

First let $24\mid h$, let $c=3w^8$ or $c=-2w^8$, with
$w\in\mathbb Z_{>0}$ and $(w,6)=1$, and put
$K=\mathbb Q(\mu_h)$ and $L=K(c^{1/h})$. We claim that the maximal
abelian subextension of $L/\mathbb Q$ is exactly $K$.
The element $c$ is a square in $K$, because $K$ contains both
$\sqrt3$ and $\sqrt{-2}$. For an odd prime $e\mid h$, an $e$th root
of $c$ generates a nonnormal real field of degree $e$: irreducibility
follows from the valuation one at three or two, and the other roots are
nonreal. Such a field cannot lie in a cyclotomic field. A fourth root
of $c$ in $K$ would similarly put $3^{1/4}$ or $2^{1/4}$ in $K$ after
division by the indicated rational square and, for $-2$, a power of
$\zeta_8$. This is again impossible. It follows that the Kummer class
of $c$ has order $h/2$. For detail, a proper prime divisor $e$ of its
possible order would give $c=z^{2e}\eta$ with
$\eta\in\mu_{h/(2e)}$; this root of unity has a $2e$th root in
$\mu_h$, contradicting the exclusions just proved.

Consequently $N=\operatorname{Gal}(L/K)$ is cyclic of order $h/2$.
Its conjugation action is cyclotomic. The subgroup $[N,G]$, where
$G=\operatorname{Gal}(L/\mathbb Q)$, is $N^2$: every cyclotomic
multiplier is odd, and the multiplier $-1$ already generates the ideal
$2$ in $\mathbb Z/(h/2)$. Thus
$N^2\subseteq[G,G]\subseteq N$. The field fixed by $N^2$ is
$K(c^{1/4})$, which contains the nonabelian splitting field of
$X^4-3$ or $X^4-2$. Hence $[G,G]$ cannot be $N^2$, and must be $N$.
This proves the claim.

If $h\mid M$ and $u\in(\mathbb Z/M)^\times$ satisfies
$u=1\pmod h$, the intersection
$L\cap\mathbb Q(\mu_M)$ is therefore $K$.
The two automorphisms $1$ on $L$ and $u$ on $\mathbb Q(\mu_M)$ agree
on that intersection, so define an automorphism of their compositum.
Chebotarev's theorem \citep[Theorem 8.31, p.~147]{MilneANT2020} supplies
infinitely many primes in the class $u\pmod M$ splitting completely
in $L$. We may exclude any finite set. This is the precise simultaneous
prime-selection statement used below.

For the local existence bound, put $d=(m,p-1)$ and assume $p\nmid mab$.
Character orthogonality on $\mathbb F_p^\times$ gives at least
\[
p-2-3(d-1)-(d-1)(d-2)\sqrt p
\]
nonzero solutions of $aX^m+bY^m=1$. Indeed the trivial character pair
contributes $p-2$, the $3(d-1)$ terms with one trivial character or
trivial product have absolute value one, and each remaining Jacobi sum
has absolute value $\sqrt p$. The latter identity follows by expanding
the two Gauss sums and changing variables $(x,y)=(zt,z(1-t))$; the
absolute value of each nontrivial Gauss sum is $\sqrt p$ by
orthogonality. If $p>(m+1)^4$, the displayed lower bound is positive:
with $T=\sqrt p$, the lower polynomial
$T^2-(m-1)(m-2)T-3m+1$ is increasing beyond $(m+1)^2$ and equals
$m^2(5m+9)>0$ there. Such a point has a unit derivative, and lifts.

We also use the following sharpened Newton estimate. If
$f\in\mathbb Z_p[T]$ and $v_p(f(1))>2e$, where
$e=v_p(f'(1))$, Newton iteration gives a root $z$ with
$v_p(z-1)\ge v_p(f(1))-e$. The first Newton increment has that
valuation, the Taylor remainder at least twice that valuation, and the
derivative keeps valuation $e$; all subsequent increments have larger
valuation. This proves both convergence and the stated bound, including
$p=2$ (compare \citealp[Theorem 7.32 and the sharper bound on p.~121]{MilneANT2020}).

Set $n_j=(j+4)!$, $B(n)=(n+1)^4$, and
$H_j(p)=2v_p(n_{j+1})+j+3$. We inductively construct
\[
a_j=3u_j^8,\quad b_j=-2v_j^8,\quad (u_jv_j,6)=1,
\quad (a_j,b_j)=1,
\]
such that $a_jX^{n_j}+b_jY^{n_j}=1$ has a unit-coordinate point at
every finite prime and
$v_p(a_j+b_j-1)\ge H_j(p)$ for every $p\le B(n_j)$.
The initial coefficients $(3,-2)$ satisfy all these conditions.

Suppose they hold at stage $j$, and write $n=n_j$, $m=n_{j+1}$,
$a=a_j$, $b=b_j$. Let $S$ contain all primes at most $B(m)$ and all
prime divisors of $ab$. Put $L_p=H_{j+1}(p)$ and
$M=\prod_{p\in S}p^{L_p}$; thus $m\mid M$ and
$L_p>2v_p(m)$. Choose old unit-coordinate solutions
$(\xi_p,\eta_p)$ at every $p\in S$. At $p\mid m$, choose them
congruent to $(1,1)$ modulo $p^{v_p(m)}$: the old error bound applies
there, and Newton in a variable with unit coefficient gives precisely
this congruence. At the other old primes any unit solution suffices.

The CRT unit class defined by the $\xi_p$ is one modulo $m$.
The prime-selection statement, applied to $\mathbb Q(\mu_n,b^{1/n})$,
gives a new prime $q$ in this class, larger than every prime of $S$
and than $n_{j+2}$. Since $b$ is an $n$th power modulo $q$,
simple-root lifting gives a unit $z_q$ with
$bz_q^n=1-aq^n$. Impose next the congruences
$r=\eta_p\pmod{p^{L_p}}$ for $p\in S$ and
$r=z_q\pmod q$. Their CRT class is still one modulo $m$.
Apply the same prime-selection statement with $h=m$ and
$c=aq^n$, which still has the required eighth-power form, to obtain a
prime $r>q$ splitting in $\mathbb Q(\mu_m,(aq^n)^{1/m})$.
Define $a_{j+1}=aq^n$, $b_{j+1}=br^n$.

At every $p\in S$ their sum is one modulo $p^{L_p}$, so Newton gives
a unit-coordinate solution of the new exponent-$m$ equation and the
next error bound. At $q$ the new equation holds at $(1,1)$ modulo $q$,
and its $Y$ derivative is a unit. At $r$ the coefficient $aq^n$ is an
$m$th power modulo $r$; setting $Y=1$ gives an equation for $X$ with a
simple nonzero root. All other primes are larger than $B(m)$ and divide
neither coefficient, so the character-sum bound applies. These are all
finite primes, completing the induction. Both prime choices can always
avoid all previous primes, and already at its first use the choice of
$q$ ranges over an infinite set.

For each fixed power modulus, the factorial exponents eventually
vanish modulo it. The coefficient products therefore converge in
$\mathcal K$ and have exactly the asserted valuation vectors. At a
fixed finite prime $p$, the valuations become constant. Subsequent
ratios are $n_j$th powers of local units; the factorials eventually
annihilate every finite quotient of $\mathbb Z_p^\times$. The two
sequences thus converge in $\mathbb Q_p^\times$ to nonzero elements
with those constant valuations. The error valuation tends to infinity,
so their sum is one. Openness of local power subgroups, also a
consequence of Newton's estimate, identifies these ordinary limits
with the localizations of the completed products. The real signs are
$+,-$, represented by the ordinary pair $(3,-2)$. Infinite support
excludes a rational pair by E.101. Taking at each stage the least
eligible residue choices and prime choices specifies an actual
terminating recursion; no numerical execution of its huge first
prime search is being asserted.
\end{proof}

\begin{result}{E.103}{Known}{SECTIONS\_OVER\_Q A.42 confirmed from rational-point finiteness}
The rational unit-equation graph
$\mathcal R=\{(\kappa(t),\kappa(1-t)):t\in\mathbb Q\setminus\{0,1\}\}$
is closed and discrete in $\mathcal K^2$. Every Cauchy net in it is
eventually constant. Thus a nonrational completed pair is separated
from all rational pairs in some simultaneous finite power quotient.
\end{result}
\begin{proof}
In a fixed fourth-power fibre choose representatives $a,b$.
Every rational pair in that fibre comes from a rational point of the
smooth projective quartic $aX^4+bY^4=Z^4$. Its genus is three, so
\citet[Satz~7, p.~365]{Faltings1983} makes this fibre finite. Each fourth-power
fibre is open and closed, and the ambient completion is Hausdorff.
Removing finitely many points gives both closedness and discreteness
of the graph. A Cauchy net eventually lies in one such finite fibre;
a common power quotient separates its finitely many points, so the
net is eventually constant. Basic open neighbourhoods are simultaneous
power-quotient fibres, proving the last assertion.
\end{proof}

\subsection{The compactness premise in Lead L2}

\begin{result}{A.101}{Refuted}{Varying ghost families are not compact; correction to Lead L2}
The original A.44 family is noncompact in $\mathcal K^2$. In particular
the set of ordinary ghosts surviving the integral class-one quotient
is noncompact. This refutes the compactness assertion in Lead L2;
it does not refute compactness of section fibres over a fixed abelian
cocycle. More generally, if the fixed-depth existence assertion of
GT\_TOWER\_CHECK C.2.12 is retained, the set of ordinary ghosts
surviving any prescribed finite integral class is noncompact.
\end{result}
\begin{proof}
For the original family, reduction modulo fifth powers gives
\[
A_5=[3q_0^{24}]_5,
\]
because $5\mid(j+4)!$ for every $j\ge1$. E.102 permits infinitely
many different $q_0$. Their valuation vectors modulo five are
different, since $24=4\pmod5$ and $q_0>3$. Thus the image in the
discrete fifth-power quotient is infinite. A compact space has finite
image in a discrete space, giving the contradiction. At class one
every pair is the abelian cocycle of a section relative to the
tangential splitting, so this example lies in the asserted survivor
set.

For the last assertion, fix a class $r$ and apply C.2.12 with
$(R,N)=(r,k!)$. It produces ghosts in the class-$r$ survivor set whose
abelian coordinates agree with $(3,-2)$ modulo $k!$. They converge in
$\mathcal K^2$ to that rational pair, which is not a ghost. Hence this
survivor set is not closed in a Hausdorff space, and cannot be compact.
This second argument is explicitly conditional on C.2.12; the first
argument requires only the arithmetic construction E.102.

The fixed-pair repair is different in kind: if the literal continuous
section fibres above one fixed abelian cocycle are compact and the
transition maps are continuous, nonemptiness at each class implies
nonemptiness of their inverse limit by the finite-intersection
property. Nothing in this argument changes, or chooses, the abelian
pair. It supplies no compactness theorem for the varying ghost locus.
\end{proof}
